% cover_letter.tex
% Cover letter for submission to Aerosol Science and Engineering (Springer)
% Compile: pdflatex manuscripts/cover_letter.tex

\documentclass[12pt,a4paper]{article}

\usepackage[utf8]{inputenc}
\usepackage[T1]{fontenc}
\usepackage{lmodern}
\usepackage[margin=2.5cm]{geometry}
\usepackage{parskip}
\usepackage{microtype}
\usepackage{hyperref}
\hypersetup{colorlinks=true, urlcolor=blue, linkcolor=blue}

\pagestyle{empty}

\begin{document}

% ── sender ────────────────────────────────────────────────────────────────────

\noindent
Federico García Crespí\\
Universidad Miguel Hernández de Elche\\
\href{mailto:fedeg@umh.es}{fedeg@umh.es}

\bigskip
\noindent
18 May 2026

\bigskip

% ── addressee ─────────────────────────────────────────────────────────────────

\noindent
Dear Editor-in-Chief,  % ← replace with name if known

\bigskip

% ── opening ───────────────────────────────────────────────────────────────────

I am pleased to submit our manuscript,

\begin{quote}
\textbf{``Does meteorology extend the useful forecast horizon for urban
PM$_{10}$?\
A multi-city, multi-horizon skill analysis with XGBoost and SARIMA''},
\end{quote}

for consideration for publication in
\textit{Aerosol Science and Engineering}.
The manuscript is original, has not been published previously, and is not
currently under consideration elsewhere.

% ── what the paper does ───────────────────────────────────────────────────────

\textbf{What we do.}
The paper addresses a practical question that is directly relevant to
operational air quality management: over how many lead hours does a
data-driven PM$_{10}$ forecasting model reliably outperform a naïve persistence
benchmark, and does the addition of concurrent meteorological covariates extend
that window?  We formalise this as the \emph{useful forecast horizon}~$H^*$,
define strict and relaxed variants, and measure it under a rigorous
rolling-origin expanding-window backtesting protocol at two European sites ---
Madrid Casa de Campo (Spain) and a network of eight stations in Ireland ---
covering a 42-month study window (January~2020 -- July~2023).

Three model classes are evaluated at each of $h = 1, \ldots, 24$ lead hours:
persistence (benchmark), SARIMA(1,0,1)(1,0,0)$_{24}$ (linear statistical
reference), and XGBoost fitted with and without nine meteorological covariates.
The Diebold--Mariano test with Harvey--Leybourne--Newbold correction is used
to assess whether observed skill differences are statistically significant.

% ── main findings ─────────────────────────────────────────────────────────────

\textbf{Main findings.}
At Madrid (lag-1 autocorrelation $\rho_1 \approx 0.90$), adding meteorology
extends $H^*_{\text{strict}}$ from 9\,h to 17\,h
($\Delta H^* = +8$\,h), a gain that is statistically significant at
$h = 12$ ($p = 0.012$, DM-HLN).  At the eight Irish stations
($\rho_1 \approx 0.82$), the mean gain is only $+0.9$\,h, with lags-only
XGBoost already achieving near-maximum horizon skill.  We show that this
contrast is mechanistically explained by the local autocorrelation-persistence
regime: when autocorrelation is high, persistence is a formidable baseline at
short leads and meteorology fills the gap that lagged-pollutant models leave
open at medium leads; when autocorrelation is moderate, lagged values already
capture the dominant variance and meteorology adds little.

% ── relevance to journal ──────────────────────────────────────────────────────

\textbf{Relevance to \textit{Aerosol Science and Engineering}.}
PM$_{10}$ is by definition an aerosol-phase quantity, and the question of
how far into the future its concentration can be reliably predicted sits
squarely within the journal's scope on aerosol measurement, modelling, and
atmospheric behaviour.  Our study provides the first systematic
multi-city analysis of the \emph{forecast horizon} for PM$_{10}$, using a
rigorous leakage-free backtesting protocol and a compact metric ($H^*$) that
translates directly into operational guidance for air quality management
systems.  The mechanistic interpretation --- linking the meteorology benefit
to the local autocorrelation-persistence regime of the aerosol time series ---
contributes new understanding of how atmospheric dynamics govern PM$_{10}$
predictability.  The revised EU Directive 2024/2881/EU, which halves the
annual PM$_{10}$ limit to 20\,\si{\micro\gram\per\cubic\metre} by 2030,
makes accurate medium-horizon (12--18\,h) forecasting increasingly urgent;
our results provide evidence-based guidance on when and where meteorological
data pipelines justify their cost in an operational context.

% ── comparison with prior rejection (AQAH) ───────────────────────────────────

\textbf{Note on prior submission.}
An earlier version of this work was submitted to \textit{Air Quality,
Atmosphere \& Health} and was not accepted following review.  The present
manuscript has been substantially revised to address the reviewers' concerns:
(i)~the mechanistic explanation linking autocorrelation regime to meteorology
benefit has been developed into a dedicated Discussion subsection; (ii)~the
multi-city design explicitly disentangles site-level predictability from model
architecture; and (iii)~both the backtesting protocol and the $H^*$ definition
are now formally specified in the Methods section.

% ── closing ───────────────────────────────────────────────────────────────────

We believe this manuscript makes a clear and timely contribution to the
literature on data-driven air quality forecasting and is well matched to your
readership.  We look forward to receiving your decision.

\bigskip

Yours sincerely,

\bigskip\bigskip

\noindent
Federico García Crespí\\
Universidad Miguel Hernández de Elche\\
\href{mailto:fedeg@umh.es}{fedeg@umh.es}

\end{document}
